Interfaces gráficas mediam a maioria das nossas interações com computadores,
seja através de \emph{laptops}, tabletes, ou navegadores web.
Desde o projeto até a implementação elas apresentam muitos desafios
\cite{myers1994}.
Além de problemas comuns de programação, como análise e processamento de dados,
destacam-se também a apresentação de dados na tela e o gerenciamento de
interações.
Para gerenciar interações na programação de interfaces é necessário coordenar
vários eventos desencadeados pelo usuário, através de dispositivos de interação
como teclado, mouse, e tela multitoque.

Para coordenar eventos é comum usar o \emph{callback}: um bloco de código executado
sempre que um dado evento ocorrer, p. ex., clique em um botão.
Esse conceito é criticado por tornar o programa complexo: a ordem de execução\footnote{A ordem de execução de um programa também é denominada \emph{fluxo
de controle}.} é
imprevisível, é definida por eventos externos, e não pela sequência especificada
pelo programador \cite{maier2010,edwards2009,fischer2007}.
“Coordenar alterações ao estado compartilhado em meio a esse caos pode ser
desconcertante, e está longe de ser modular. A definição coloquial é \emph{Inferno de
Callbacks}.”\footnote{Tradução literal do inglês \emph{‘Callback Hell’}: expressão
popular na comunidade de programação.} \cite[p. 2; tradução nossa]{edwards2009}.
Inerentemente imperativo, o \emph{callback} é muito usado na \emph{programação orientada a
objetos} (POO) através do \emph{Observer Pattern}\footnote{O \emph{Observer Pattern} é usado para coordenar eventos em
linguagens orientadas a objetos, e \emph{callback} às vezes é chamado de \emph{event
handler}, \emph{event listener} ou \emph{observer}, mas em essência o conceito é o mesmo.}
\cite{blackheath2016,maier2010}.
Para esclarecer os desafios enfrentados por sistemas de software em produção,
podemos citar uma análise das aplicações \emph{desktop} da Adobe, de 2005, na qual se
concluiu que a lógica de coordenação de eventos consistia em um terço do
código, além de conter metade dos \emph{bugs} reportados durante o ciclo de vida do
produto \cite{jarvi2008}.

Uma alternativa é a \emph{programação reativa} (PR), proposta para aplicações
orientadas a eventos \cite{bainomugisha2013} e aplicada a interfaces gráficas,
animações, aplicações web, robótica e redes de sensores \cite{salvaneschi2017}.
Na PR, o programador declara as relações entre os valores, e o ambiente de
execução faz todas as atualizações que elas exigem \cite{salvaneschi2017}.
Esse comportamento é o das planilhas eletrônicas, como o \emph{Google Sheets} e o
\emph{Microsoft Excel}, “possivelmente a linguagem de programação mais utilizada
por usuários finais” \cite[p. 2; tradução nossa]{bainomugisha2013}.
A PR descende da \emph{programação funcional} (PF): a maior parte da pesquisa na
área vem do Fran, uma linguagem funcional do fim dos anos 1990
\cite[p. 2]{bainomugisha2013}.

Argumenta-se há tempos que a PR torna os programas mais fáceis de compreender
que o \emph{Observer Pattern}, mas faltava evidência empírica
\cite{salvaneschi2017}.
Um experimento controlado comparou a compreensão de programas escritos com PR
e com o \emph{Observer Pattern}, primeiro com 38 participantes
\cite{salvaneschi2014} e depois com mais 89, num total de 127
\cite{salvaneschi2017}.
Os participantes eram estudantes do quarto ano de Ciência da Computação,
divididos ao acaso em dois grupos.
Um grupo leu dez programas em Scala escritos com PR, e o outro, os mesmos
programas escritos com o \emph{Observer Pattern}.
Para cada programa, os participantes respondiam a uma pergunta de múltipla
escolha sobre o comportamento dele.
O grupo da PR acertou mais, com diferença significativa (médias de 8,13 e 7,05
acertos em dez), e não levou mais tempo para responder.
Além disso, no grupo do \emph{Observer Pattern} os acertos acompanharam o nível de
habilidade de programação dos participantes, o que não se observou no grupo da
PR \cite{salvaneschi2017}.

Para \textcite[p. 1]{moseley2006}, o que mais contribui para a complexidade
dos sistemas grandes é o tratamento do estado, ao lado do volume de código e
da preocupação explícita com o fluxo de controle\footnote{\textcite{moseley2006} distinguem a complexidade essencial,
própria do problema, da acidental, todo o resto, e sustentam que a maior parte
da que resta nos sistemas é acidental.}.
Eles propõem reduzi-la com uma abordagem baseada em programação funcional e no
modelo relacional de dados.
Conceitos de programação declarativa podem, assim, mitigar problemas do
desenvolvimento de software em larga escala.

Na programação de interfaces gráficas, a coordenação entre evento, estado e
tela se escreve, nas tecnologias mais usadas, em três notações, e duas delas
são declarativas.
Na \emph{imperativa com callbacks}, cada \emph{callback} altera o estado e a tela
diretamente, como no jQuery \cite{openjs2026} e nos Web Components sem
biblioteca \cite{mozilla2026}.
Nas declarativas, o programador descreve a tela em função do estado, e o
ambiente de execução a mantém atualizada \cite{meta2026a}.
Na \emph{declarativa por re-renderização}, cada mudança de estado reexecuta o
componente, como no React \cite{meta2026}.
Na \emph{reativa fina com signals}, cada \emph{signal} avisa os valores e trechos da
tela que dependem dele, e só esses são recalculados, como no Solid e no
Angular com \emph{signals} \cite{solidjs2026,google2026}.
A notação reativa fina aplica conceitos da PR, como os valores que variam no
tempo, chamados \emph{signals}, e a propagação automática de alterações
\cite{salvaneschi2015}.

Cada tecnologia fica na notação em que escreve a maior parte da
coordenação, mas a classificação não é exclusiva.
Em quase todo problema, cada parte pede o seu conjunto de conceitos
\cite[p. 10]{roy2009}, e as implementações deste trabalho também os
misturam.
O jQuery e o Web Component derivam a Lista filtrável com \texttt{filter}, \texttt{sort} e
\texttt{map}, funções de PF.
E o React e o Solid recorrem a efeitos para sincronizar o componente com
sistemas externos, como o armazenamento do navegador no Carrinho.
A documentação do React chama os efeitos de “uma saída de emergência do
paradigma do React” \cite[tradução nossa]{meta2026d}.

As pesquisas de uso medem tecnologias, não notações, e por isso o uso de
cada notação se afere pelas tecnologias que a adotam.
No Stack Overflow Developer Survey 2025, React, jQuery e Angular estão entre
as sete tecnologias web mais usadas pelos respondentes, com 44,7\%, 23,4\% e
18,2\% \cite{stackoverflow2025}.
O jQuery está ainda em 65,5\% dos \emph{sites} que o W3Techs monitora
\cite{w3techs2026}.
Nenhuma das duas listas traz o Solid ou os Web Components.
O Solid é usado por 10\% dos respondentes do State of JS 2025
\cite{devographics2026}.
Por isso, o uso da notação reativa fina se apoia no Angular, que também a
adota, embora o número do Angular não separe quem usa \emph{signals}.
Da mesma forma, o uso da notação imperativa se apoia no jQuery.
O Web Component entra por ser o modelo de componente do próprio navegador,
sem biblioteca \cite{mozilla2026}, e nele a notação é o \texttt{addEventListener}
cujo \emph{callback} altera a tela.
Essas medidas têm ainda limites de amostra.
O Stack Overflow recruta os respondentes sobretudo pelos próprios canais, o
que favorece os usuários mais ativos do site \cite{stackoverflow2025}.
O State of JS é aberto a quem quiser responder e se declara o retrato de um
subconjunto de desenvolvedores \cite{devographics2026}.
O W3Techs conta uma biblioteca como usada por um \emph{site} se a achar em
qualquer página dele \cite{w3techs2026a}, o que favorece o jQuery e nada diz
das aplicações novas.
Pergunta-se, então: como as notações mais usadas na prática para programar
interfaces gráficas, a imperativa com \emph{callbacks}, a declarativa por
re-renderização e a reativa fina com \emph{signals}, se comparam quanto à
usabilidade, segundo as Dimensões Cognitivas de Notações?

O objetivo geral é comparar a usabilidade das três notações segundo as
Dimensões Cognitivas de Notações, na programação de interfaces típicas da
web, em TypeScript, com replicação no Android, em Kotlin. Especificamente:

\begin{itemize}
\item Demonstrar, com processamento de listas, os conceitos de PF em que se
apoiam as notações declarativas;
\item Implementar cinco casos de interface (Contador, Formulário com validação,
Busca com sugestões, Lista filtrável e Carrinho) com Web Component,
jQuery, React, Solid e Angular com \emph{signals}, e o Contador, o Formulário e
a Lista no Android, com Views e Jetpack Compose;
\item Avaliar as implementações por oito Dimensões Cognitivas;
\item Sintetizar as vantagens e desvantagens de cada notação por problema de
coordenação, como estado derivado e assincronia.
\end{itemize}

Esta pesquisa é de natureza \emph{aplicada}, e quanto aos objetivos ela se
classifica como \emph{exploratória}, pois tem o “\textelp{} foco mais aberto para
investigação de fenômenos (culturais, sociais, técnicos, históricos, etc.)
pouco sistematizados e/ou passíveis de várias perspectivas de interpretação.”
\cite[p. 32]{leal2011}.
Quanto aos meios empregados este trabalho constitui um \emph{estudo de casos
múltiplos}:

\begin{citacao}
  O estudo de casos múltiplos – denominado, em algumas áreas, como
  administração pública e ciência política, de método de caso comparativo – é
  preferido quando há possibilidade de comparar semelhanças e de contrastar
  diferenças entre os casos selecionados. \cite[p. 43]{leal2011}
\end{citacao}

\textcite{yin2001} salienta que pesquisas desse tipo envolvem o estudo profundo
e minucioso de um ou mais casos, que neste trabalho são programas concretos
implementados para demonstrar os conceitos de programação.
A linguagem \emph{JavaScript} é usada na implementação dos casos.
Essa escolha é resguardada pelo fato de que, além de ser a língua franca da
web, ela tem suporte para os paradigmas estudados.

Primeiro, o paradigma de PF é demonstrado com a implementação de programas
para \emph{processamento de listas}.
Isso fundamenta diferenças essenciais entre conceitos declarativos e
imperativos.
Em seguida, conceitos de PR e o \emph{callback} são demonstrados na \emph{coordenação
de eventos} em interfaces gráficas.
Essa segunda parte esclarece conceitos declarativos do paradigma de PR e o
tradicional \emph{callback}, imperativo.

A fim de contrastar as vantagens e desvantagens de cada conceito, a linguagem de
programação em cada caso é analisada através de um conjunto de critérios
padronizados, chamados de \emph{Dimensões Cognitivas de Notações} (DCs), do inglês
\emph{Cognitive Dimensions of Notations} \cite{green1989}.
Criadas para analisar a usabilidade de ‘artefatos de
informação’\footnote{Geralmente sistemas de software, especialmente linguagens de
programação. Mais informações podem ser encontradas no site
\url{http://www.cl.cam.ac.uk/\~afb21/CognitiveDimensions/}.}, essas DCs já foram usadas para investigar API
— \emph{Application Programming Interface} \cite{clarke2003}, recursos de linguagens
de programação \cite{sadowski2011}, e paradigmas de programação \cite{kiss2014}.
